% ============================================================
% tests/test-commandes.tex
% Fichier de test de non-régression pour outilsprof
% ============================================================
% Compilation : xelatex test-commandes.tex
% ============================================================

\documentclass[12pt]{extarticle}

% --- Polices (XeLaTeX) ---
\usepackage{fontspec}
\setmainfont{Times New Roman}
\setsansfont{Helvetica}
\setmonofont{Menlo}

% --- Maths et tableaux ---
\usepackage{amsmath, amssymb}
\usepackage{array}
\usepackage{booktabs}
\usepackage{enumitem}
\usepackage{fancyhdr}
\usepackage{multirow}
\usepackage{multicol}

% --- Graphiques et boîtes ---
\usepackage{tikz}
\usepackage{geometry}                % sans options (tcolorbox l'inclut)
\usepackage{tcolorbox}
\tcbuselibrary{skins,breakable}

% --- Le package à tester ---
\usepackage{outilsprof}

% --- Réglages globaux ---
\geometry{top=2cm, bottom=2cm, left=1cm, right=1cm}
\setlength{\headheight}{14pt}
\setlength{\parindent}{0pt}
\setlength{\parskip}{1em}

% --- Notation pour la case grise attendue ---
\newcommand{\grise}{\ensuremath{\square}}

\begin{document}

\begin{center}
  {\Large\bfseries Tests de non-régression — \texttt{outilsprof} v1.11.0}
\end{center}

% ============================================================
\section*{1. Masquage \texttt{|} — additions}
% ============================================================

\subsection*{1.1 — Le 4 est masqué}

\textbf{Entrée :} \verb!\addition[solution]{1|4, 56}!
\quad\textbf{Attendu :} $1\grise + 56 = 70$

\addition[solution]{1|4, 56}

\subsection*{1.2 — Le 1 est masqué}

\textbf{Entrée :} \verb!\addition[solution]{|14, 56}!
\quad\textbf{Attendu :} $\grise 4 + 56 = 70$

\addition[solution]{|14, 56}

\subsection*{1.3 — Un 0 implicite est masqué (fin de nombre)}

\textbf{Entrée :} \verb!\addition[solution]{47|, 12}!
\quad\textbf{Attendu :} $47\grise + 12 = 482$

\addition[solution]{47|, 12}

\subsection*{1.4 — Plusieurs masques}

\textbf{Entrée :} \verb!\addition[solution]{6|34, |72|8}!
\quad\textbf{Attendu :} $6\grise 4 + \grise 2\grise = 1362$

\addition[solution]{6|34, |72|8}

% ============================================================
\section*{2. Masquage \texttt{|} — soustractions}
% ============================================================

\subsection*{2.1 — \texttt{8|3, 47|}}

\soustraction[solution]{8|3, 47|}

\subsection*{2.2 — \texttt{|73|4, 2|78}}

\soustraction[solution]{|73|4, 2|78}

\subsection*{2.1b — Vérification du 0 implicite : \texttt{8|3, 47|}}

\textbf{Calcul attendu :} $470 - 83 = 387$
\quad(Les termes sont permutés : 83 < 470)

\soustraction[solution]{8|3, 47|}

% ============================================================
\section*{3. Masquage \texttt{|} — multiplications}
% ============================================================

\subsection*{3.1 — \texttt{24, 5|6}}

\multiplication[solution]{24, 5|6}

\subsection*{3.2 — \texttt{|1|2|3, |45}}

\multiplication[solution]{|1|2|3, |45}

% ============================================================
\section*{4. Masquage \texttt{|} — durées}
% ============================================================

\subsection*{4.1 — \texttt{3+1|2, 5+30}}

\additiondurees[solution]{3+1|2, 5+30}

\subsection*{4.2 — \texttt{3+|5+45, 2+30}}

\additiondurees[solution]{3+|5+45, 2+30}

% ============================================================
\section*{5. Masquage \texttt{|} — angles}
% ============================================================

\subsection*{5.1 — \texttt{5|6+30, 12+45}}

\additionangles[solution]{5|6+30, 12+45}

% ============================================================
\section*{6. Division euclidienne et décimale}
% ============================================================

\subsection*{6.1 — Division euclidienne simple}

\division[dividende=1234, diviseur=56, solution]

\subsection*{6.2 — Division euclidienne avec étapes}

\division[dividende=1234, diviseur=56, solution, etapes={112, 112}]

\subsection*{6.3 — Division décimale : $22 \div 7$}

\division[dividende=22, diviseur=7, solution, decimales=3,
          etapes={21, 7, 28, 14}]

\subsection*{6.4 — Division décimale : $1000 \div 70$}

\division[dividende=1000, diviseur=70, solution, decimales=5,
          etapes={70, 280, 140, 560, 350, 490, 70}]

\subsection*{6.5 — Padding à gauche : $1 \div 8$}

\division[dividende=1, diviseur=8, solution, decimales=3]

\subsection*{6.6 — Padding à droite : $12 \div 8$}

\division[dividende=12, diviseur=8, solution, decimales=2]

% ============================================================
\section*{7. Autres commandes (test de fumée)}
% ============================================================

\subsection*{7.1 — Axe gradué}

\axegradue[fin=400, pas=100, nombre=5,
           points={A/120, B/340/blue, C/60/green}]

\subsection*{7.2 — Repère gradué}

\reperegradue[bg={(-5,-5)}, hd={(5,5)},
              points={A/(4,5), B/(-2,3)/blue},
              affichernom=true, affichercoords=true]

\subsection*{7.3 — Boîtes}

\begin{retenir}
  Le théorème de Pythagore : $a^2 + b^2 = c^2$.
\end{retenir}

\begin{vigilance}
  Ne pas oublier de vérifier les unités.
\end{vigilance}

\begin{methode}
  1. Identifier les données \\
  2. Appliquer la formule
\end{methode}

\subsection*{7.4 — Grandeurs, angles, durées}

\unite{9,81}{m/s^2} \quad
\degre{56+30+15} \quad
\dureeheure{5+45+30}

\subsection*{7.5 — Lignes de réponse}

Le résultat est \reponse. \grandeReponse[0.5]

\end{document}